% !TeX root = ../main.tex

\documentclass[../main.tex]{subfiles}
\begin{document}

\chapter{KẾT LUẬN VÀ KIẾN NGHỊ}

\section{Kết luận}

Đồ án tốt nghiệp với đề tài điều chỉnh và thuyết minh công nghệ dự án "Nhà máy xử lý chất thải rắn sinh hoạt thành phố Hội An" đã giải quyết một cách toàn diện và khoa học bài toán quản lý chất thải đang vô cùng cấp bách tại địa phương. Qua quá trình nghiên cứu, tính toán và phân tích đa chiều, đồ án rút ra những kết luận quan trọng sau:

\textbf{Thứ nhất}, việc chuyển đổi cấu hình từ công nghệ "khí hóa -- tuabin khí" sang "nhiệt phân -- chưng cất -- phát điện diesel" là một quyết định điều chỉnh có tính thực tiễn và chiến lược cao. Rác thải sinh hoạt Hội An mang đặc tính độ ẩm cao và độ đồng nhất thấp, khiến cho phương án khí hóa gặp nhiều rào cản về mặt xử lý kỹ thuật và rủi ro ngưng trệ. Công nghệ nhiệt phân khắc phục hoàn toàn điểm yếu này nhờ khả năng tiếp nhận nguyên liệu linh hoạt hơn và sản xuất ra sản phẩm chủ đạo là dầu sinh học -- một dạng năng lượng dễ dàng lưu trữ, pha trộn và ứng dụng.

\textbf{Thứ hai}, quy trình công nghệ được thiết lập mang đậm tính kinh tế tuần hoàn. Việc khép kín các chu trình vật chất và năng lượng bên trong nhà máy: từ việc dùng khí syngas duy trì nhiệt độ lò nhiệt phân, tái sử dụng phần nhiệt xả từ khói thải động cơ để sấy khô rác đầu vào, cho đến việc cố định hóa tro bay và tuần hoàn nước thải 100\% đã minh chứng cho khả năng vận hành tự cấp tự túc, giảm mức tiêu hao năng lượng từ bên ngoài xuống mức cực thấp.

\textbf{Thứ ba}, qua các kết quả tính toán cân bằng năng lượng, dự án không chỉ hoàn thành sứ mệnh xử lý triệt để 120 tấn rác sinh hoạt mỗi ngày mà còn đóng vai trò như một trạm phát điện năng lượng tái tạo quy mô nhỏ. Với sản lượng điện tinh cung cấp lên lưới đạt 1,82 MW, nhà máy đóng góp xấp xỉ 16 GWh điện năng hàng năm, trực tiếp cắt giảm hàng ngàn tấn khí thải nhà kính so với kịch bản nhiệt điện hóa thạch truyền thống.

\textbf{Thứ tư}, các đánh giá về mặt tác động môi trường khẳng định rằng, dưới sự vận hành nghiêm ngặt của hệ thống lọc bụi tĩnh điện ESP, tháp hấp thụ đa cấp, cùng trạm xử lý nước rỉ rác công nghệ màng MBR, nhà máy cam kết không phát thải thứ cấp độc hại ra môi trường sống xung quanh. Đây là nền tảng cốt lõi đảm bảo sự đồng thuận của cộng đồng dân cư và chính quyền sở tại.

Nhìn chung, giải pháp công nghệ nhiệt phân kết hợp phát điện này đáp ứng đầy đủ cả ba tiêu chí phát triển bền vững: Hiệu quả kinh tế - An toàn xã hội - Thân thiện môi trường.

\section{Kiến nghị}

Để dự án được triển khai thành công và đưa vào vận hành thương mại với hiệu suất tối ưu, đồ án đề xuất một số kiến nghị hướng tới các bên liên quan:

\textbf{Đối với Chủ đầu tư và Ban quản lý dự án:}
\begin{itemize}
    \item Cần khẩn trương thiết lập bộ máy nhân sự kỹ thuật nòng cốt ngay từ giai đoạn lắp đặt và chạy thử nghiệm thiết bị. Việc để nhân sự trực tiếp làm việc cùng các chuyên gia chuyển giao công nghệ sẽ giúp quá trình làm chủ hệ thống vận hành về sau diễn ra trơn tru.
    \item Chú trọng xây dựng quy trình SOP (Standard Operating Procedure) chuẩn hóa cho mọi công đoạn, đặc biệt là khâu pha trộn tỷ lệ dầu sinh học với dầu diesel thương mại để bảo vệ độ bền động cơ phát điện.
    \item Thiết lập quỹ dự phòng rủi ro môi trường, đảm bảo luôn sẵn sàng các vật tư hóa chất, than hoạt tính, dung dịch kiềm để xử lý sự cố trong mọi tình huống.
\end{itemize}

\textbf{Đối với Chính quyền địa phương (Thành phố Hội An và tỉnh Quảng Nam):}
\begin{itemize}
    \item Thúc đẩy mạnh mẽ công tác truyền thông và triển khai chương trình phân loại rác tại nguồn triệt để đối với các hộ gia đình, cơ sở kinh doanh, đặc biệt là hệ thống nhà hàng, khách sạn. Nếu rác được tách riêng rác hữu cơ, rác tái chế và rác vô cơ ngay từ đầu, công đoạn tiền xử lý của nhà máy sẽ giảm tải đáng kể, chất lượng viên nén RDF sẽ tăng cao, kéo theo hiệu suất phát điện tăng vọt.
    \item Cần có chính sách hỗ trợ giá mua điện ưu đãi (Feed-in Tariff - FIT) đối với các dự án điện rác tương tự nhằm tạo động lực mạnh mẽ thu hút các nhà đầu tư tư nhân tham gia vào lĩnh vực xử lý môi trường.
    \item Hỗ trợ tạo kênh tiêu thụ đầu ra cho sản phẩm than sinh học biochar thông qua việc giới thiệu, kết nối với các mô hình nông nghiệp công nghệ cao, nông nghiệp sinh thái trên địa bàn tỉnh.
\end{itemize}

\end{document}